% Sommario in italiano (stesura del 2026-11-17). Numeri dai capitoli 3-6; aggiornarli se cambieranno.
Le event camera registrano i cambiamenti di luminosità nel momento in cui avvengono, e producono un segnale in cui l'informazione sta soprattutto nell'ordine temporale dell'attività. Gli Spiking Transformer sono candidati naturali per elaborarlo, ma trattano il tempo in modo in gran parte implicito, affidandolo al potenziale di membrana dei neuroni, e rendere esplicito l'accesso al passato costa memoria, una risorsa scarsa nei modelli compatti pensati per le FPGA. Questa tesi si chiede come spendere un budget di circa mezzo milione di parametri perché uno Spiking Transformer colga le relazioni temporali che contano nei dati a eventi, tenendo esplicito il costo della memoria.

Il lavoro modifica un solo fattore alla volta, su quattro dataset neuromorfici, e descrive il costo con conteggi analitici di parametri, operazioni e stato persistente. Riallocare il budget di una versione compatta di QKFormer, con un frontend piramidale, un mixer causale a più ritardi e una convoluzione depthwise, porta la Macro-F1 di DVS-Lip da 44 a 55 punti con appena 320 parametri in più. Una rappresentazione dell'ingresso a lunghezza variabile, in cui ogni passo ha la stessa durata fisica, dà ai ritardi della rete un significato temporale preciso e migliora nettamente DVS128 Gesture e THU-EACT-50-CHL. Su questa base la tesi propone una memoria relazionale a intervalli, l'IRM antisimmetrico, che riassume il passato in bande causali a distanze fisiche fisse e lo confronta con il presente attraverso un termine bilineare a basso rango, ispirato al correlatore di movimento di Hassenstein e Reichardt. Su THU-EACT-50-CHL l'IRM migliora il modello di riferimento di $7{,}6$ punti di Macro-F1 sulla validazione e di $10{,}8$ sul test ufficiale, un guadagno che i controlli di pari capacità non riproducono, e una volta compresso, senza riaddestramento, costa meno del mixer che sostituisce. Su DVS-Lip il beneficio è piccolo, e il limite del modello compatto si rivela legato alla generalizzazione, che capacità e distillazione riescono invece a spostare.

I risultati, ottenuti per lo più con un solo addestramento per configurazione e con costi stimati anziché misurati su hardware, indicano che in uno Spiking Transformer compatto la capacità temporale dipende meno dalla quantità di memoria che dalla scelta di che cosa sia un passo, di quale passato conservare e di come confrontarlo con il presente.
